% OR01-S001 | lapisan asli: tujuan, batas, dan notasi
% segment-id: d90.orig.v1.tr01.seg0001
\chapter{Metode Stokastik Komposit, Cermin, dan Minibatch}
\label{orig01:chapter:stochastic-composite}

Bab ini mengisi penghubung yang belum tersedia di tulang punggung Habring
dan tiga modul Becker yang telah diterima. Habring telah mengembangkan
operator proksimal, gradien proksimal deterministik, dan subgradien stokastik
terproyeksi. Modul Becker~3 telah membuktikan ketakbiasan penaksir SAGA dan
identitas varians tabelnya. Yang masih diperlukan ialah satu kerangka yang
menjelaskan bagaimana operator proksimal, geometri non-Euklides, pengambilan
sampel minibatch, dan reduksi varians saling terhubung.

Seluruh rumusan, bukti penghubung, latihan, petunjuk, solusi, dan laboratorium
dalam bab ini ditulis secara mandiri. Catatan Habring, catatan gradien
stokastik Cl\'ement W. Royer, modul Becker~3, serta rujukan penelitian pada
akhir bab dipakai untuk memeriksa istilah, batas teori, dan konteks; tidak ada
prosa, tata letak, gambar, atau kode mereka yang disalin.

Kita memakai norma $\norm{\cdot}$ pada ruang berdimensi hingga dan norma
dual
\begin{equation}
  \norm{y}_*=\sup_{\norm{x}\leq1}\inner{y}{x}.
\end{equation}
Untuk bagian proksimal dan minibatch, norma tersebut adalah norma Euklides.
Filtrasi $\mathcal F_k$ memuat seluruh informasi sebelum oracle pada iterasi
$k$ dipanggil. Dengan demikian, $x_k$ terukur terhadap $\mathcal F_k$.

\section{Masalah komposit dan oracle stokastik}
\label{orig01:section:model}

Pertimbangkan masalah komposit
\begin{equation}
  \label{orig01:eq:composite-problem}
  \min_{x\in\R^d} F(x),
  \qquad F(x)=f(x)+g(x),
\end{equation}
dengan $f:\R^d\to\R$ konveks, terdiferensialkan, dan mempunyai gradien
$L$-Lipschitz, sedangkan $g:\R^d\to(-\infty,+\infty]$ proper, konveks, dan
semikontinu bawah. Andaikan himpunan peminim $X^*=\arg\min F$ tidak kosong.

Oracle pada $x_k$ mengembalikan $G_k$ dan memenuhi, hampir pasti,
\begin{equation}
  \label{orig01:eq:oracle}
  \E[G_k\mid\mathcal F_k]=\nabla f(x_k),
  \qquad
  \E\!\left[\norm{G_k-\nabla f(x_k)}^2\mid\mathcal F_k\right]
  \leq \nu_k^2.
\end{equation}
Besaran $\nu_k^2$ boleh berubah sepanjang iterasi. Notasi ini sengaja
memisahkan varians oracle dari konstanta kemulusan $L$ dan dari fungsi
kerugian $\ell$ pada model statistik.

% OR01-S002 | lapisan asli: gradien proksimal stokastik
% segment-id: d90.orig.v1.tr01.seg0002
\section{Gradien proksimal stokastik}
\label{orig01:section:spg}

Dengan ukuran langkah $\tau_k>0$, pembaruan gradien proksimal stokastik ialah
\begin{equation}
  \label{orig01:eq:spg-update}
  x_{k+1}
  =\prox_{\tau_k g}\bigl(x_k-\tau_kG_k\bigr).
\end{equation}
Jika $g=\delta_C$ adalah fungsi indikator himpunan konveks tertutup $C$,
pembaruan ini menjadi proyeksi stokastik. Jika $g(x)=\lambda\norm{x}_1$,
pembaruan ini menjadi ambang lunak setelah langkah gradien stokastik.

\begin{quote}
\textbf{Algoritma 1: gradien proksimal stokastik.}
Pilih $x_0\in\dom g$. Untuk $k=0,1,\dots$, panggil oracle pada $x_k$ untuk
memperoleh $G_k$, pilih $\tau_k$, lalu hitung~\eqref{orig01:eq:spg-update}.
Untuk pelaporan ergodik, simpan rerata iterat yang disebutkan dalam teorema
berikut.
\end{quote}

\begin{lemma}[Ketaksamaan satu langkah]
  \label{orig01:lemma:spg-one-step}
  Ambil $x\in\dom g$ dan tetapkan $\tau_k=\tau$ dengan
  $0<\tau\leq(2L)^{-1}$. Di bawah~\eqref{orig01:eq:oracle}, iterasi
  \eqref{orig01:eq:spg-update} memenuhi
  \begin{multline}
    \label{orig01:eq:spg-one-step}
    \E[F(x_{k+1})-F(x)\mid\mathcal F_k]\
    \leq
    \frac{\norm{x_k-x}^2-
                \E[\norm{x_{k+1}-x}^2\mid\mathcal F_k]}{2\tau}
    +\tau\nu_k^2.
  \end{multline}
\end{lemma}

\begin{proof}
  Tuliskan $x^+=x_{k+1}$,
  $s=x^+-x_k$, dan $\delta_k=G_k-\nabla f(x_k)$. Kondisi optimalitas
  proks memberikan
  \begin{equation}
    \frac{x_k-x^+}{\tau}-G_k\in\partial g(x^+).
  \end{equation}
  Ketaksamaan subgradien untuk $g$, kekonveksan $f$, dan lemma penurunan untuk
  $f$ yang $L$-mulus menghasilkan
  \begin{multline}
    F(x^+)-F(x)
    \leq
    \frac{\norm{x_k-x}^2-\norm{x^+-x}^2}{2\tau}
    -\frac{1-L\tau}{2\tau}\norm{s}^2
    -\inner{\delta_k}{x^+-x}.
    \label{orig01:eq:spg-pathwise}
  \end{multline}
  Pecah suku terakhir sebagai
  $-\inner{\delta_k}{s}-\inner{\delta_k}{x_k-x}$. Karena
  $x_k-x$ terukur terhadap $\mathcal F_k$ dan
  $\E[\delta_k\mid\mathcal F_k]=0$, suku kedua hilang dalam ekspektasi
  bersyarat. Untuk suku pertama, maksimisasi kuadrat atau ketaksamaan Young
  memberi
  \begin{equation}
    -\inner{\delta_k}{s}
    -\frac{1-L\tau}{2\tau}\norm{s}^2
    \leq\frac{\tau}{2(1-L\tau)}\norm{\delta_k}^2
    \leq\tau\norm{\delta_k}^2,
  \end{equation}
  dengan ketaksamaan terakhir menggunakan $\tau L\leq1/2$. Mengambil
  ekspektasi bersyarat membuktikan klaim.
\end{proof}

\begin{theorem}[Batas ergodik proksimal stokastik]
  \label{orig01:theorem:spg-ergodic}
  Andaikan $\nu_k^2\leq\sigma^2/b$ untuk semua $k$, dengan $b\geq1$.
  Ambil $x^*\in X^*$, $0<\tau\leq(2L)^{-1}$, dan definisikan
  \begin{equation}
    \bar x_K=\frac1K\sum_{k=0}^{K-1}x_{k+1}.
  \end{equation}
  Maka, untuk $K\geq1$,
  \begin{equation}
    \label{orig01:eq:spg-ergodic-rate}
    \E[F(\bar x_K)-F(x^*)]
    \leq
    \frac{\norm{x_0-x^*}^2}{2\tau K}
    +\frac{\tau\sigma^2}{b}.
  \end{equation}
\end{theorem}

\begin{proof}
  Terapkan Lema~\ref{orig01:lemma:spg-one-step} dengan $x=x^*$, ambil
  ekspektasi penuh, lalu jumlahkan untuk $k=0,\dots,K-1$. Suku jarak
  meneleskop dan suku terakhir yang tidak negatif dapat dibuang. Kekonveksan
  $F$ memberi
  $F(\bar x_K)\leq K^{-1}\sum_{k=0}^{K-1}F(x_{k+1})$, yang menyelesaikan
  bukti.
\end{proof}

Jika $R=\norm{x_0-x^*}>0$, $\sigma>0$, dan batas kemulusan tidak aktif,
pilihan yang
menyeimbangkan kedua suku pada~\eqref{orig01:eq:spg-ergodic-rate} ialah
\begin{equation}
  \tau=R\sqrt{\frac{b}{2\sigma^2K}},
  \qquad
  \E[F(\bar x_K)-F(x^*)]
  \leq\frac{\sqrt2R\sigma}{\sqrt{bK}}.
\end{equation}
Dalam praktik, $R$ dan $\sigma$ biasanya tidak diketahui. Rumus ini adalah
penjelas skala, bukan resep penalaan tanpa diagnosis.

% OR01-S003 | lapisan asli: minibatch
% segment-id: d90.orig.v1.tr01.seg0003
\section{Minibatch: identitas varians dan biaya oracle}
\label{orig01:section:minibatch}

Ambil $b$ sampel bersyarat iid $G_{k,1},\dots,G_{k,b}$ dengan rerata
$\nabla f(x_k)$ dan varians kuadrat paling besar $\sigma^2$. Penaksir
minibatch dengan penggantian adalah
\begin{equation}
  \label{orig01:eq:minibatch}
  \widehat G_k^{(b)}=\frac1b\sum_{j=1}^{b}G_{k,j}.
\end{equation}

\begin{prop}[Penyusutan varians dengan penggantian]
  \label{orig01:prop:minibatch-with-replacement}
  Penaksir~\eqref{orig01:eq:minibatch} tak bias dan memenuhi
  \begin{equation}
    \E\!\left[
      \norm{\widehat G_k^{(b)}-\nabla f(x_k)}^2
      \mid\mathcal F_k
    \right]\leq\frac{\sigma^2}{b}.
  \end{equation}
\end{prop}

\begin{proof}
  Setelah dikondisikan pada $\mathcal F_k$, simpangan setiap sampel mempunyai
  rerata nol dan pasangan simpangan yang berbeda saling bebas. Semua suku
  silang dalam kuadrat norma mempunyai ekspektasi nol. Tersisa $b$ suku
  diagonal yang masing-masing dikalikan $b^{-2}$.
\end{proof}

Untuk jumlah hingga $f=N^{-1}\sum_{i=1}^{N}f_i$, pengambilan sampel tanpa
penggantian mempunyai koreksi populasi hingga yang tidak boleh disamakan
dengan rumus iid.

\begin{prop}[Koreksi populasi hingga]
  \label{orig01:prop:finite-population}
  Andaikan $N\geq2$ dan $1\leq b\leq N$. Pada suatu $x$, tuliskan
  $h_i=\nabla f_i(x)-\nabla f(x)$ dan
  $V_N=N^{-1}\sum_{i=1}^{N}\norm{h_i}^2$. Jika $S$ dipilih seragam dari semua
  subhimpunan berukuran $b$ tanpa penggantian, maka
  \begin{equation}
    \label{orig01:eq:finite-population}
    \E\norm{\frac1b\sum_{i\in S}\nabla f_i(x)-\nabla f(x)}^2
    =\frac{N-b}{b(N-1)}V_N.
  \end{equation}
\end{prop}

\begin{proof}
  Peluang inklusi satu indeks adalah $b/N$ dan peluang inklusi dua indeks
  berbeda adalah $b(b-1)/(N(N-1))$. Kembangkan kuadrat norma. Karena
  $\sum_i h_i=0$, berlaku
  $\sum_{i\ne j}\inner{h_i}{h_j}=-\sum_i\norm{h_i}^2$. Substitusi kedua
  peluang inklusi lalu penyederhanaan memberikan faktor pada
  \eqref{orig01:eq:finite-population}. Khusus $b=N$, varians tepat nol.
\end{proof}

Minibatch menurunkan varians per iterasi, tetapi memakai $b$ evaluasi gradien
komponen. Jika biaya serial yang tersedia adalah $T=bK$, batas tertala ideal
$\sqrt2R\sigma/\sqrt{bK}$ menjadi $\sqrt2R\sigma/\sqrt T$: tidak ada
percepatan kompleksitas oracle hanya dari pengelompokan. Keuntungan nyata
dapat datang dari paralelisme, arsitektur perangkat, pengurangan komunikasi,
atau batas langkah yang lebih baik. Pernyataan ini juga menjelaskan mengapa
perbandingan hanya berdasarkan banyak iterasi dapat menyesatkan.

% OR01-S004 | lapisan asli: geometri cermin
% segment-id: d90.orig.v1.tr01.seg0004
\section{Penurunan cermin stokastik}
\label{orig01:section:mirror}

Proyeksi Euklides memperlakukan semua arah dengan geometri yang sama. Pada
simpleks, matriks semidefinit positif, atau ruang dengan struktur sparsitas,
pilihan geometri lain dapat lebih alami.

Misalkan $X\subset\R^d$ tak kosong, tertutup, dan konveks. Ambil fungsi
pembangkit $h:X\to\R$ yang bernilai hingga dan kontinu pada $X$,
terdiferensialkan di setiap iterat, dan $\alpha$-konveks kuat terhadap
$\norm{\cdot}$. Untuk setiap $x\in X$ dan setiap titik keterdiferensialan $y$,
\begin{equation}
    h(x)\geq h(y)+\inner{\nabla h(y)}{x-y}
  +\frac\alpha2\norm{x-y}^2.
\end{equation}
Divergensi Bregman yang dibangkitkan oleh $h$ ialah
\begin{equation}
  \label{orig01:eq:bregman}
  D_h(x,y)=h(x)-h(y)-\inner{\nabla h(y)}{x-y}.
\end{equation}
Divergensi ini tidak harus simetris dan bukan metrik, tetapi
$D_h(x,y)\geq(\alpha/2)\norm{x-y}^2$.

Untuk fungsi konveks $\Phi:X\to\R$, andaikan oracle menghasilkan $G_k$
dengan
\begin{equation}
  \label{orig01:eq:mirror-oracle}
  \bar G_k\coloneqq\E[G_k\mid\mathcal F_k]\in\partial\Phi(x_k),
  \qquad
  \E[\norm{G_k}_*^2\mid\mathcal F_k]\leq M^2.
\end{equation}
Ambil $\tau_k>0$ dan andaikan masalah berikut hampir pasti mencapai peminim
unik $x_{k+1}$ pada titik keterdiferensialan $h$:
\begin{equation}
  \label{orig01:eq:mirror-update}
  x_{k+1}=\arg\min_{x\in X}
  \left\{\tau_k\inner{G_k}{x}+D_h(x,x_k)\right\}.
\end{equation}
Untuk pembanding optimum $x^*$ yang dipakai di bawah, andaikan pula
$D_h(x^*,x_0)<\infty$. Asumsi eksplisit ini dapat diganti dengan syarat
Legendre/interior standar yang menjamin sifat-sifat yang sama.

\begin{quote}
\textbf{Algoritma 2: penurunan cermin stokastik.}
Pilih $x_0$ dalam daerah terdiferensialkan $h$. Pada iterasi $k$, ambil
subgradien stokastik $G_k$, selesaikan masalah Bregman
\eqref{orig01:eq:mirror-update}, dan perbarui rerata berbobot
$\sum_j\tau_jx_j/\sum_j\tau_j$.
\end{quote}

\begin{lemma}[Ketaksamaan tiga titik stokastik]
  \label{orig01:lemma:mirror-one-step}
  Untuk setiap $x\in X$, pembaruan~\eqref{orig01:eq:mirror-update} memenuhi
  \begin{equation}
    \label{orig01:eq:mirror-one-step}
    \tau_k\inner{G_k}{x_k-x}
    \leq D_h(x,x_k)-D_h(x,x_{k+1})
    +\frac{\tau_k^2}{2\alpha}\norm{G_k}_*^2.
  \end{equation}
\end{lemma}

\begin{proof}
  Kondisi optimalitas pada himpunan $X$ memberi
  \begin{equation}
    \inner{\tau_kG_k+\nabla h(x_{k+1})-\nabla h(x_k)}
           {x-x_{k+1}}\geq0.
  \end{equation}
  Identitas tiga titik Bregman mengubah bagian yang memuat $h$ menjadi
  \begin{equation}
    \tau_k\inner{G_k}{x_{k+1}-x}
    \leq D_h(x,x_k)-D_h(x,x_{k+1})-D_h(x_{k+1},x_k).
  \end{equation}
  Tambahkan $\tau_k\inner{G_k}{x_k-x_{k+1}}$ pada kedua sisi. Gunakan
  $D_h(x_{k+1},x_k)\geq(\alpha/2)\norm{x_{k+1}-x_k}^2$, dualitas norma, dan
  ketaksamaan Young. Hasilnya tepat~\eqref{orig01:eq:mirror-one-step}.
\end{proof}

\begin{theorem}[Batas ergodik cermin stokastik]
  \label{orig01:theorem:mirror-ergodic}
  Di bawah~\eqref{orig01:eq:mirror-oracle}, untuk $K\geq1$ definisikan
  $S_K=\sum_{k=0}^{K-1}\tau_k$ dan
  \begin{equation}
    \widetilde x_K=\frac1{S_K}\sum_{k=0}^{K-1}\tau_kx_k.
  \end{equation}
  Untuk setiap $x^*\in\arg\min_{x\in X}\Phi(x)$,
  \begin{equation}
    \label{orig01:eq:mirror-rate}
    \E[\Phi(\widetilde x_K)-\Phi(x^*)]
    \leq
    \frac{D_h(x^*,x_0)}{S_K}
    +\frac{M^2}{2\alpha}
      \frac{\sum_{k=0}^{K-1}\tau_k^2}{S_K}.
  \end{equation}
\end{theorem}

\begin{proof}
  Ambil ekspektasi bersyarat pada
  Lema~\ref{orig01:lemma:mirror-one-step}. Karena $x_k-x^*$ terukur terhadap
  $\mathcal F_k$ dan $\bar G_k\in\partial\Phi(x_k)$,
  \begin{equation}
    \E[\inner{G_k}{x_k-x^*}\mid\mathcal F_k]
    \geq\Phi(x_k)-\Phi(x^*).
  \end{equation}
  Jumlahkan sehingga divergensi Bregman meneleskop, lalu gunakan kekonveksan
  $\Phi$ pada rerata berbobot.
\end{proof}

Untuk $M>0$ dan $D_h(x^*,x_0)>0$, gunakan
$\tau=\sqrt{2\alpha D_h(x^*,x_0)/(M^2K)}$. Ruas kanan kemudian menjadi
\begin{equation}
  M\sqrt{\frac{2D_h(x^*,x_0)}{\alpha K}}.
\end{equation}
Pada geometri Euklides, $h(x)=\norm{x}_2^2/2$ dan pembaruan cermin menjadi
subgradien stokastik terproyeksi. Jadi metode cermin memperluas, bukan
menduplikasi, teorema stokastik Habring.

\subsection{Simpleks dan pembaruan eksponensial}
\label{orig01:subsection:entropy}

Pada simpleks
$\Delta_d=\{x\in\R_+^d:\sum_{j=1}^{d}x_j=1\}$, ambil
$h(x)=\sum_jx_j\log x_j$, dengan konvensi perluasan kontinu
$0\log0=0$. Fungsi ini kontinu pada $\Delta_d$ dan terdiferensialkan pada
interior relatifnya. Mulai dari $x_{0,j}>0$. Divergensi Bregman ialah
divergensi Kullback--Leibler. Terhadap norma $\ell_1$, fungsi ini
$1$-konveks kuat. Kondisi optimalitas memberi pembaruan tertutup
\begin{equation}
  \label{orig01:eq:exponentiated-update}
  x_{k+1,j}
  =\frac{x_{k,j}\exp(-\tau_kG_{k,j})}
         {\sum_{r=1}^{d}x_{k,r}\exp(-\tau_kG_{k,r})}.
\end{equation}
Jika $x_0$ seragam, maka $D_h(x^*,x_0)\leq\log d$. Dengan batas
$\norm{G_k}_\infty\leq M$, Teorema~\ref{orig01:theorem:mirror-ergodic}
memberi skala $M\sqrt{2\log(d)/K}$ setelah penalaan ideal. Ketergantungan
logaritmik pada dimensi menjelaskan kegunaan geometri entropi di simpleks.

% OR01-S005 | lapisan asli: penghubung reduksi varians
% segment-id: d90.orig.v1.tr01.seg0005
\section{Dari SAGA ke pembaruan proksimal}
\label{orig01:section:prox-saga}

Untuk masalah jumlah hingga
$f=N^{-1}\sum_i f_i$, modul Becker~3 mendefinisikan penaksir SAGA
\begin{equation}
  v_k=\nabla f_{J_k}(x_k)-g_{J_k}^k+\bar g_k
\end{equation}
dan membuktikan
$\E[v_k\mid\mathcal F_k]=\nabla f(x_k)$. Untuk masalah komposit, perubahan
algoritmik yang benar bukan mengurangkan subgradien $g$, melainkan memakai
langkah proksimal
\begin{equation}
  \label{orig01:eq:prox-saga}
  x_{k+1}=\prox_{\tau g}(x_k-\tau v_k).
\end{equation}

\begin{cor}[Penghubung varians untuk Prox-SAGA]
  \label{orig01:cor:prox-saga-bridge}
  Andaikan asumsi masalah komposit berlaku, $0<\tau\leq(2L)^{-1}$, dan
  definisikan
  \begin{equation}
    \omega_k^2
    =\E[\norm{v_k-\nabla f(x_k)}^2\mid\mathcal F_k].
  \end{equation}
  Untuk $\bar x_K=K^{-1}\sum_{k=0}^{K-1}x_{k+1}$ dan $x^*\in X^*$,
  \begin{equation}
    \label{orig01:eq:prox-saga-bridge}
    \E[F(\bar x_K)-F(x^*)]
    \leq\frac{\norm{x_0-x^*}^2}{2\tau K}
    +\frac{\tau}{K}\sum_{k=0}^{K-1}\E[\omega_k^2].
  \end{equation}
  Jika $\E[\omega_k^2]\leq C\rho^k$ dengan $C\geq0$ dan $0<\rho<1$, maka
  \begin{equation}
    \E[F(\bar x_K)-F(x^*)]
    \leq\frac{\norm{x_0-x^*}^2}{2\tau K}
    +\frac{\tau C}{K(1-\rho)}.
  \end{equation}
\end{cor}

\begin{proof}
  Terapkan bukti Teorema~\ref{orig01:theorem:spg-ergodic} dengan
  $G_k=v_k$ dan varians bersyarat $\omega_k^2$. Untuk klaim terakhir, gunakan
  jumlah deret geometri hingga yang dibatasi oleh $(1-\rho)^{-1}$.
\end{proof}

Korolari ini menjelaskan hubungan yang hilang, tetapi sengaja tidak mengklaim
bahwa varians SAGA selalu turun geometrik. Pembuktian klaim semacam itu
memerlukan fungsi Lyapunov gabungan untuk galat iterat dan galat tabel.
Teorema laju linear bersyarat pada modul Becker~3 menyediakan hasil khusus
untuk kasus mulus dan konveks kuat; korolari di sini menunjukkan bagaimana
penaksir yang sama masuk ke masalah komposit.

% OR01-S006 | lapisan asli: laboratorium
% segment-id: d90.orig.v1.tr01.seg0006
\section{Laboratorium 1: regresi renggang dengan biaya oracle tetap}
\label{orig01:section:lab}

Laboratorium terbuka pada
\texttt{labs/original-01/stochastic-composite-lab.py} membangun masalah
regresi kuadrat dengan regularisasi $\ell_1$,
\begin{equation}
  \label{orig01:eq:lab-objective}
  F(x)=\frac{1}{2N}\norm{Ax-y}_2^2+\lambda\norm{x}_1.
\end{equation}
Skrip membandingkan gradien proksimal stokastik, minibatch dengan penggantian,
dan Prox-SAGA pada anggaran evaluasi gradien komponen yang sama. Referensi
numerik dihitung oleh FISTA deterministik sampai norma pemetaan gradien
proksimal memenuhi toleransi yang dilaporkan.

Tugas laboratorium:
\begin{enumerate}
  \item jalankan skrip tanpa mengubah benih acak dan verifikasikan konfigurasi
    beku yang dicatat pada berkas hasil;
  \item bandingkan kesenjangan objektif terhadap evaluasi gradien komponen,
    bukan hanya terhadap iterasi;
  \item periksa apakah minibatch menurunkan varians arah pada checkpoint yang
    sama;
  \item ubah ukuran minibatch tetapi pertahankan anggaran oracle, lalu
    jelaskan perbedaan antara keuntungan statistik dan keuntungan paralel;
  \item ganti $\lambda$ dan ukur perubahan sparsitas serta norma pemetaan
    gradien proksimal;
  \item laporkan setiap hasil yang bertentangan dengan batas teori beserta
    asumsi yang mungkin tidak terpenuhi.
\end{enumerate}
CSV dan JSON yang dihasilkan adalah permukaan aksesibel utama. Grafik SVG
hanya merupakan bantuan visual dan tidak menjadi satu-satunya pembawa
informasi.

Dengan konfigurasi beku ($N=320$, $d=40$, $\lambda=0{,}03$, 12 epoch, dan
3.840 evaluasi gradien komponen per metode), nilai referensi FISTA adalah
$0{,}31517396778742246$ dengan norma pemetaan gradien proksimal
$1{,}49\times10^{-15}$. Hasil terminal deterministik adalah:
\begin{center}
\begin{tabular}{lrrr}
\hline
Metode & Kesenjangan objektif & Norma pemetaan & Tak nol\\
\hline
Proks-SGD, $b=1$ & $4{,}9169\times10^{-2}$ & $2{,}1155\times10^{-2}$ & 33\\
Proks-minibatch, $b=16$ & $2{,}5707\times10^{-3}$ & $9{,}8630\times10^{-3}$ & 10\\
Prox-SAGA & $2{,}0444\times10^{-9}$ & $9{,}8666\times10^{-6}$ & 5\\
\hline
\end{tabular}
\end{center}
Angka tersebut adalah satu eksperimen terkontrol, bukan urutan kinerja
universal. Berkas hasil JSON, CSV, dan SVG mempunyai SHA-256 berikut:
\begin{center}\small\ttfamily
86ff701a\ldots c447\\
61a6591a\ldots 5d37\\
87c772d9\ldots 2830.
\end{center}
Identitas lengkap dicatat dalam receipt QA.

% OR01-S007 | lapisan asli: latihan, petunjuk, solusi
% segment-id: d90.orig.v1.tr01.seg0007
\section{Latihan, petunjuk, dan solusi lengkap}
\label{orig01:section:exercises}

\begin{exercise}[Pembaruan proksimal untuk Lasso]
  \label{orig01:exercise:lasso-prox}
  Untuk $f(x)=(2N)^{-1}\norm{Ax-y}_2^2$ dan
  $g(x)=\lambda\norm{x}_1$, tuliskan pembaruan
  \eqref{orig01:eq:spg-update} ketika minibatch $S_k$ dipakai. Tunjukkan
  bahwa setiap koordinat diperoleh dengan ambang lunak.

  \textbf{Petunjuk bertahap.}
  (i) Hitung gradien setiap suku kuadrat.
  (ii) Gunakan keterpisahan norma $\ell_1$.
  (iii) Selesaikan masalah proksimal skalar pada tiga kasus tanda.

  \textbf{Solusi lengkap.}
  Jika $f_i(x)=\tfrac12(a_i^\top x-y_i)^2$, maka
  \begin{equation}
    G_k=\frac1{|S_k|}\sum_{i\in S_k}a_i(a_i^\top x_k-y_i).
  \end{equation}
  Tuliskan $z_k=x_k-\tau_kG_k$. Karena
  $\prox_{\tau_k\lambda\norm{\cdot}_1}$ terpisah per koordinat,
  \begin{equation}
    (x_{k+1})_j
    =\operatorname{sign}((z_k)_j)
      \max\{|(z_k)_j|-\tau_k\lambda,0\}.
  \end{equation}
  Rumus ini mengikuti dengan memeriksa kondisi
  $0\in u-z+\tau_k\lambda\partial|u|$ pada $u>0$, $u<0$, dan $u=0$.
\end{exercise}

\begin{exercise}[Konstanta pada ketaksamaan satu langkah]
  \label{orig01:exercise:young-constant}
  Mulai dari~\eqref{orig01:eq:spg-pathwise}. Buktikan batas yang lebih tajam
  \begin{equation}
    \E[F(x_{k+1})-F(x)\mid\mathcal F_k]
    \leq\frac{\norm{x_k-x}^2-
      \E\norm{x_{k+1}-x}^2}{2\tau}
    +\frac{\tau\nu_k^2}{2(1-L\tau)}
  \end{equation}
  untuk $0<\tau<L^{-1}$.

  \textbf{Petunjuk bertahap.}
  (i) Maksimalkan $-\inner{\delta}{s}-c\norm{s}^2$ terhadap $s$.
  (ii) Ambil $c=(1-L\tau)/(2\tau)$.
  (iii) Gunakan ketakbiasan hanya setelah memisahkan $x_k-x$.

  \textbf{Solusi lengkap.}
  Dualitas norma Euklides dan pelengkapan kuadrat memberikan
  \begin{equation}
    \sup_s\{-\inner{\delta}{s}-c\norm{s}^2\}
    =\frac{\norm{\delta}^2}{4c}
    =\frac{\tau\norm{\delta}^2}{2(1-L\tau)}.
  \end{equation}
  Suku $-\inner{\delta_k}{x_k-x}$ mempunyai ekspektasi bersyarat nol.
  Substitusi batas varians menyelesaikan klaim. Lema utama memakai
  $\tau L\leq1/2$ hanya untuk mengganti faktor ini dengan batas sederhana
  $\tau\nu_k^2$.
\end{exercise}

\begin{exercise}[Sampel tanpa penggantian]
  \label{orig01:exercise:finite-population}
  Untuk $N\geq2$, buktikan Proposisi~\ref{orig01:prop:finite-population} dengan peubah
  indikator $I_i=\1\{i\in S\}$. Periksa kasus $b=1$ dan $b=N$.

  \textbf{Petunjuk bertahap.}
  (i) Gunakan $\E I_i=b/N$.
  (ii) Untuk $i\ne j$, gunakan
  $\E[I_iI_j]=b(b-1)/(N(N-1))$.
  (iii) Pakai $\sum_i h_i=0$.

  \textbf{Solusi lengkap.}
  Karena $b^{-1}\sum_{i\in S}h_i=b^{-1}\sum_iI_ih_i$,
  ekspektasi kuadrat normanya adalah
  \begin{multline}
    \frac1{b^2}\left[
      \frac bN\sum_i\norm{h_i}^2
      +\frac{b(b-1)}{N(N-1)}
        \sum_{i\ne j}\inner{h_i}{h_j}
    \right].
  \end{multline}
  Suku silang sama dengan $-\sum_i\norm{h_i}^2$. Koefisien akhirnya
  $(N-b)/(bN(N-1))$, dan penggantian
  $\sum_i\norm{h_i}^2=NV_N$ memberi rumus yang diminta. Untuk $b=1$ faktor
  menjadi satu; untuk $b=N$ menjadi nol.
\end{exercise}

\begin{exercise}[Cermin entropi pada simpleks]
  \label{orig01:exercise:entropy-update}
  Turunkan~\eqref{orig01:eq:exponentiated-update} dari
  \eqref{orig01:eq:mirror-update}. Jelaskan mengapa iterat yang mulai positif
  tetap positif dan berjumlah satu.

  \textbf{Petunjuk bertahap.}
  (i) Tambahkan pengali Lagrange untuk $\sum_jx_j=1$.
  (ii) Diferensialkan terhadap setiap $x_j$.
  (iii) Gunakan kendala jumlah untuk menentukan faktor normalisasi.

  \textbf{Solusi lengkap.}
  Dengan konvensi $0\log0=0$ dan
  $D_h(x,x_k)=\sum_jx_j\log(x_j/x_{k,j})$ pada simpleks, kondisi
  stasioner adalah
  \begin{equation}
    \tau_kG_{k,j}+\log(x_j/x_{k,j})+1+\eta=0.
  \end{equation}
  Jadi $x_j=x_{k,j}\exp(-\tau_kG_{k,j})\exp(-1-\eta)$. Menjumlahkan seluruh
  koordinat menunjukkan bahwa faktor terakhir adalah kebalikan jumlah pada
  penyebut~\eqref{orig01:eq:exponentiated-update}. Eksponensial dan penyebut
  positif mempertahankan kepositifan, sedangkan normalisasi mempertahankan
  jumlah satu.
\end{exercise}

\begin{exercise}[Anggaran oracle dan ukuran minibatch]
  \label{orig01:exercise:oracle-budget}
  Andaikan batas tertala ideal adalah
  $\sqrt2R\sigma/\sqrt{bK}$ dan satu iterasi minibatch memakai $b$ evaluasi
  gradien komponen. Untuk anggaran serial $T=bK$, tentukan ketergantungannya
  pada $b$. Sebutkan dua keadaan ketika minibatch tetap berguna.

  \textbf{Petunjuk bertahap.}
  (i) Ganti $K$ dengan $T/b$.
  (ii) Bedakan waktu dinding dari jumlah panggilan oracle.
  (iii) Perhatikan pengambilan sampel tanpa penggantian.

  \textbf{Solusi lengkap.}
  Substitusi memberi
  \begin{equation}
    \frac{\sqrt2R\sigma}{\sqrt{b(T/b)}}
    =\frac{\sqrt2R\sigma}{\sqrt T},
  \end{equation}
  yang bebas dari $b$. Jadi model ideal tidak memprediksi penghematan
  evaluasi gradien serial. Minibatch tetap berguna jika evaluasi dapat
  diparalelkan atau komunikasi dapat diamortisasi; ia juga dapat membantu
  melalui koreksi populasi hingga tanpa penggantian, vektorisasi perangkat,
  atau batas langkah yang tidak tercakup dalam model sederhana.
\end{exercise}

\begin{exercise}[Varians yang menyusut pada Prox-SAGA]
  \label{orig01:exercise:prox-saga-decay}
  Mulai dari~\eqref{orig01:eq:prox-saga-bridge}. Jika
  $\E\omega_k^2\leq C/(k+1)$, turunkan batas eksplisit menggunakan bilangan
  harmonik $H_K=\sum_{j=1}^{K}j^{-1}$. Bandingkan dengan kasus geometrik.

  \textbf{Petunjuk bertahap.}
  (i) Jumlahkan batas varians.
  (ii) Gunakan $H_K\leq1+\log K$.
  (iii) Jangan menyimpulkan laju iterat terakhir dari batas ergodik.

  \textbf{Solusi lengkap.}
  Substitusi langsung memberi
  \begin{equation}
    \E[F(\bar x_K)-F(x^*)]
    \leq\frac{\norm{x_0-x^*}^2}{2\tau K}
    +\frac{\tau C H_K}{K}
    \leq\frac{\norm{x_0-x^*}^2}{2\tau K}
    +\frac{\tau C(1+\log K)}{K}.
  \end{equation}
  Peluruhan geometrik menjadikan jumlah varians terbatas secara seragam dan
  menghasilkan $\Oc(1/K)$ tanpa faktor logaritmik. Kedua kesimpulan hanya
  berlaku untuk rerata $\bar x_K$ dalam korolari ini; laju iterat terakhir
  memerlukan analisis tambahan.
\end{exercise}

% OR01-S008 | lapisan asli: peta asumsi dan rujukan
% segment-id: d90.orig.v1.tr01.seg0008
\section{Peta asumsi dan batas klaim}
\label{orig01:section:assumption-map}

\begin{center}
\begin{tabular}{p{0.22\textwidth}p{0.32\textwidth}p{0.35\textwidth}}
\hline
Hasil & Asumsi penentu & Yang tidak diklaim\\
\hline
Proksimal stokastik & $f$ konveks dan $L$-mulus; $g$ proper, tertutup,
konveks; oracle tak bias; varians terbatas & laju iterat terakhir atau
konvergensi nonkonveks\\
Minibatch iid & sampel bersyarat bebas dengan penggantian & koreksi populasi
hingga\\
Minibatch tanpa penggantian & subhimpunan seragam berukuran
$1\leq b\leq N$ dari populasi dengan $N\geq2$ &
independensi antaranggota batch\\
Cermin stokastik & pembangkit $\alpha$-konveks kuat; pembaruan terdefinisi;
$D_h(x^*,x_0)<\infty$; momen dual kedua terbatas & simetri divergensi
Bregman\\
Prox-SAGA penghubung & ketakbiasan SAGA dan kontrol varians yang dinyatakan &
peluruhan geometrik varians tanpa Lyapunov\\
\hline
\end{tabular}
\end{center}

Bab ini tidak membahas ketaksamaan variasional atau operator monoton maksimal;
materi tersebut adalah tranche asli berikutnya. Bab ini juga tidak mengulang
LP/MIP, simpleks sebagai algoritma pemrograman linear, dualitas LP, sensitivitas
LP, jaringan, atau optimisasi diskret yang merupakan cakupan O018. Kata
\emph{simpleks} di sini hanya menunjuk himpunan probabilitas $\Delta_d$.

\section*{Rujukan matematis dan saksi verifikasi}
\begingroup\sloppy
\begin{itemize}
  \item Andreas Habring, \emph{Lecture Notes: Convex Optimization},
    arXiv:2607.11664v1, khususnya operator proksimal dan subgradien stokastik.
  \item Cl\'ement W. Royer, \emph{Lecture Notes on Stochastic Gradient
    Methods}, edisi 2023/2024, khususnya pembahasan minibatch dan biaya epoch;
    dipakai sebagai saksi perbandingan, bukan sumber prosa bab ini.
  \item Lorenzo Rosasco, Silvia Villa, dan Bang C\^ong V\~u,
    \emph{Convergence of Stochastic Proximal Gradient Algorithm},
    arXiv:1403.5074, sebagai rujukan primer bagi keluarga metode proksimal
    stokastik.
  \item Amir Beck dan Marc Teboulle, \emph{Mirror Descent and Nonlinear
    Projected Subgradient Methods for Convex Optimization},
    \emph{Operations Research Letters} 31 (2003), 167--175,
    DOI:10.1016/S0167-6377(02)00231-6.
  \item Aaron Defazio, Francis Bach, dan Simon Lacoste-Julien,
    \emph{SAGA: A Fast Incremental Gradient Method With Support for
    Non-Strongly Convex Composite Objectives}, arXiv:1407.0202v3.
\end{itemize}
\endgroup
